%=============================================================================
% Section 10: Future Directions
%=============================================================================
\section{Future Directions}

\subsection{Near-Term Priorities (6 Months)}

\textbf{Capability Discovery Protocol.} Future work could define a standardized capability discovery protocol enabling dynamic capability registration, capability health monitoring, and automatic capability indexing.

\textbf{Status.} The ai-lib project has initiated work on this protocol, with draft specifications for:
\begin{itemize}
    \item Provider heartbeat and health reporting
    \item Capability deprecation notices
    \item Version negotiation for capability schemas
\end{itemize}

\textbf{Benchmark Suite.} A cross-runtime benchmark suite is planned to evaluate:
\begin{itemize}
    \item Translation overhead across Rust, Python, TypeScript, and Go runtimes
    \item Throughput under concurrent synchronous and streaming workloads
    \item Memory footprint and startup latency comparisons
\end{itemize}

\subsection{Medium-Term Priorities (12--18 Months)}

\textbf{Application Packaging Standard.} A standardized application packaging format would enable:
\begin{itemize}
    \item Portable AI application distribution
    \item Standard deployment models
    \item Capability requirement declaration
\end{itemize}

\textbf{Design Considerations:}
\begin{itemize}
    \item Manifest-based capability requirements (similar to package.json or Cargo.toml)
    \item Container-native packaging for cloud deployment
    \item Version-pinned capability declarations for reproducibility
\end{itemize}

\textbf{Edge Runtime Support.} The execution layer architecture naturally supports edge computing scenarios. Runtimes can operate locally on:
\begin{itemize}
    \item Personal devices (smartphones, tablets)
    \item Home servers (NAS devices, personal cloud)
    \item Embedded systems (IoT gateways, automotive systems)
\end{itemize}

This architecture enables personal AI infrastructure, where users operate local runtimes orchestrating personal capabilities.

\subsection{Long-Term Vision (18+ Months)}

\textbf{Capability Marketplace.} As capabilities become standardized, marketplaces may emerge for:
\begin{itemize}
    \item Premium capability offerings
    \item Capability quality ratings and reviews
    \item Usage-based pricing and metering
\end{itemize}

\textbf{Asynchronous Notification Support.} While the current design relies on polling for long-running tasks, future runtime versions may support callback-based notifications. In this model:
\begin{itemize}
    \item Providers register webhook endpoints in their manifests
    \item The runtime translates provider-specific callbacks into protocol-standard events
    \item The Contact Layer subscribes to capability completion events rather than polling
\end{itemize}

This would reduce connection overhead for high-concurrency scenarios and enable true event-driven capability orchestration.

\textbf{Stateless Constraint Clarification.} The runtime itself remains stateless: it acts purely as a protocol translator and forwarder, pushing standardized events to the Contact Layer, which maintains all subscription state and routing policy.

\subsection{Research Agenda}

Beyond engineering priorities, several research questions remain open:
\begin{enumerate}
    \item \textbf{Formal Verification.} Can capability contracts be formally verified for provider compliance?
    \item \textbf{Economic Modeling.} How do pricing models influence routing strategies at the Contact Layer?
    \item \textbf{Security Boundaries.} What is the minimal trusted computing base for capability execution?
\end{enumerate}

These questions extend the architectural formalization into formal methods, computational economics, and security domains.
